\documentclass[12pt]{article}
\usepackage[a4paper,margin=1in]{geometry}
\usepackage[utf8]{inputenc}
\usepackage[T1]{fontenc}
\usepackage{amsmath,amssymb}
\usepackage{textcomp}
\usepackage{parskip}
\begin{document}
\begin{center}
{\Large \textbf{2017 STEP II}}
\end{center}

\subsection*{Page 1}
\begin{flushleft}
\bigskip\textbf{Section A: Pure Mathematics}\par
d 1\par
\[1 Note: In this question you may use without proof the result arctan x = .\]
\[dx 1 + x^2\]
\[\leq \geq\]
Let\par
n\par
\[In = x arctan x dx ,\]
(\par
\[where n = 0, 1, 2, 3, . . . .\]
\textbf{(i)} Show that, for n \leq 0 ,\par
\[\pi 1 x n+1\]
\[(n + 1)In = dx\]
\[4 - ( 1 + x^2\]
and evaluate I .\par
\textbf{(ii)} Find an expression, in terms of n, for (n + 3)In+2 + (n + 1)In .\par
Use this result to evaluate I .\par
\textbf{(iii)} Prove by induction that, for n \leq 1,\par
2n\par
1 1\par
r\par
\[(4n + 1)I 4n = A ( 1) ,\]
- 2 - r\par
\[)r=1\]
where A is a constant to be determined.\par
\bigskip\textbf{Question 2.} The sequence of numbers x , x , x , . . . satisfies\par
0 1 2\par
axn 1\par
\[xn+1 = - .\]
xn + b\par
(You may assume that a, b and x\par
\[are such that xn + b = 0 .)\]
\[\leq\]
Find an expression for xn+2 in terms of a, b and xn.\par
\textbf{(i)} Show that a + b = 0 is a necessary condition for the sequence to be periodic with\par
period 2.\par
\[Note: The sequence is said to be periodic with period k if xn+k = xn for all n, and\]
\[there is no integer m with 0 < m < k such that xn+m = xn for all n.\]
\textbf{(ii)} Find necessary and sufficient conditions for the sequence to have period 4.\par
\[\geq\geq\theta\mu\nu) \phi \geq\]
\end{flushleft}
\subsection*{Page 2}
\begin{flushleft}
\bigskip\textbf{Question 3.} (i) Sketch, on x-y axes, the set of all points satisfying sin y = sin x, for \pi \geq x \geq \pi and\par
-\par
\[\pi \geq y \geq \pi. You should give the equations of all the lines on your sketch.\]
-\par
\textbf{(ii)} Given that\par
\[sin y = 1 sin x\]
obtain an expression, in terms of x, for y \leq when 0 \geq x \geq 1 \pi and 0 \geq y \geq 1 \pi, and show\par
2 2\par
that\par
3 sin x\par
\[\leq\leq\]
\[y = .\]
\[- (4 sin2 x) 3 2\]
-\par
Use these results to sketch the set of all points satisfying sin y = 1 sin x for 0 \geq x \geq 1 \pi\par
2 2\par
\[and 0 \geq y \geq 1 \pi.\]
\[Hence sketch the set of all points satisfying sin y = 1 sin x for \pi \geq x \geq \pi and\]
2 -\par
\[\pi \geq y \geq \pi .\]
-\par
\textbf{(iii)} Without further calculation, sketch the set of all points satisfying cos y = 1 sin x for\par
\[\pi \geq x \geq \pi and \pi \geq y \geq \pi.\]
- -\par
\bigskip\textbf{Question 4.} The Schwarz inequality is\par
b 2 b b\par
\[f(x) g(x) dx \geq f(x) 2 dx g(x) 2 dx . ( )\]
\[\int(a \sum \int(a \sum \int(a \sum ∗\]
\[\leq \geq \leq \geq\]
\textbf{(i)} By setting f(x) = 1 in ( ), and choosing g(x), a and b suitably, show that for t > 0 ,\par
∗\par
t\par
e 1 t\par
\[- \geq .\]
et + 1 2\par
\textbf{(ii)} By setting f(x) = x in ( ), and choosing g(x) suitably, show that\par
∗\par
\[e - 2 1 x^2 dx \leq 12 1 e - 4 1 )2 .\]
( -\par
\[0 \leq\]
\textbf{(iii)} Use ( ) to show that\par
\[∗ 1 \pi\]
\[64 2 \pi\]
\[\geq \sqrt\{sin\} x dx \geq .\]
\[25\pi ( ( 2\]
\[\geq\geq\theta\mu\nu) \pi\]
\end{flushleft}
\subsection*{Page 3}
\begin{flushleft}
\bigskip\textbf{Question 5.} A curve C is determined by the parametric equations\par
\[x = at2 , y = 2at ,\]
\[where a > 0 .\]
\textbf{(i)} Show that the normal to C at a point P , with non-zero parameter p, meets C again\par
at a point N , with parameter n, where\par
\[n = p + .\]
\[- \int p \sum\]
\textbf{(ii)} Show that the distance P N is given by\par
| |\par
\[(p^2 + 1)3\]
\[P N 2 = 16a2\]
\[| | p^4\]
\[and that this is minimised when p^2 = 2 .\]
\textbf{(iii)} The point Q, with parameter q, is the point at which the circle with diameter P N\par
cuts C again. By considering the gradients of QP and QN , show that\par
2q\par
\[2 = p^2 q^2 + .\]
- p\par
Deduce that P N is at its minimum when Q is at the origin.\par
| |\par
6 Let\par
n\par
\[Sn = ,\]
\[\sqrt\{r\}\]
\[)r=1\]
where n is a positive integer.\par
\textbf{(i)} Prove by induction that\par
\[Sn \geq 2\sqrt\{n\} 1 .\]
-\par
\textbf{(ii)} Show that (4k + 1)\sqrtk + 1 > (4k + 3)\sqrtk for k \leq 0 .\par
Determine the smallest number C such that\par
\[Sn \leq 2\sqrt\{n\} + C .\]
\[2\sqrt\{n\} -\]
\[\geq\geq\theta\mu\nu) \sigma \geq\]
\end{flushleft}
\subsection*{Page 4}
\begin{flushleft}
\bigskip\textbf{Question 7.} The functions f and g are defined, for x > 0, by\par
\[f(x) = x\]
x\par
\[, g(x) =\]
xf(x)\par
.\par
\textbf{(i)} By taking logarithms, or otherwise, show that f(x) > x for 0 < x < 1 . Show further\par
\[that x < g(x) < f(x) for 0 < x < 1 .\]
\[Write down the corresponding results for x > 1 .\]
\[\leq\]
\textbf{(ii)} Find the value of x for which f (x) = 0 .\par
\textbf{(iii)} Use the result x ln x 0 as x 0 to find lim f(x), and write down lim g(x) .\par
\[\to \to x\to0 x\to0\]
\textbf{(iv)} Show that x -1 + ln x \leq 1 for x > 0.\par
\[\leq\]
\[Using this result, or otherwise, show that g (x) > 0 .\]
\[Sketch the graphs, for x > 0, of y = x, y = f(x) and y = g(x) on the same axes.\]
8 All vectors in this question lie in the same plane.\par
The vertices of the non-right-angled triangle ABC have position vectors a, b and c, respec-\par
tively. The non-zero vectors u and v are perpendicular to BC and CA, respectively.\par
Write down the vector equation of the line through A perpendicular to BC, in terms of u, a\par
and a parameter λ.\par
The line through A perpendicular to BC intersects the line through B perpendicular to CA\par
at P . Find the position vector of P in terms of a, b, c and u.\par
Hence show that the line CP is perpendicular to the line AB.\par
\[\geq\geq\theta\mu\nu) \omega\]
\end{flushleft}
\subsection*{Page 5}
\begin{flushleft}
\bigskip\textbf{Section B: Mechanics}\par
\bigskip\textbf{Question 9.} Two identical rough cylinders of radius r and weight W rest, not touching each other but a\par
negligible distance apart, on a horizontal floor. A thin flat rough plank of width 2a, where\par
a < r, and weight kW rests symmetrically and horizontally on the cylinders, with its length\par
parallel to the axes of the cylinders and its faces horizontal. A vertical cross-section is shown\par
in the diagram below.\par
The coefficient of friction at all four contacts is 1 . The system is in equilibrium.\par
\textbf{(i)} Let F be the frictional force between one cylinder and the floor, and let R be the\par
normal reaction between the plank and one cylinder. Show that\par
\[R sin θ = F (1 + cos θ) ,\]
where θ is the acute angle between the plank and the tangent to the cylinder at the\par
point of contact.\par
\[Deduce that 2 sin θ \geq 1 + cos θ .\]
\textbf{(ii)} Show that\par
2 1 + cos θ\par
\[N = 1 + F ,\]
\[\int k \sum \int sin θ \sum\]
where N is the normal reaction between the floor and one cylinder.\par
Write down the condition that the cylinder does not slip on the floor and show that it\par
is satisfied with no extra restrictions on θ.\par
\textbf{(iii)} Show that sin θ \geq 4 and hence that r \geq 5a .\par
\[\geq\geq\theta\mu\nu) ) \geq\]
\end{flushleft}
\subsection*{Page 6}
\begin{flushleft}
\bigskip\textbf{Question 10.} A car of mass m makes a journey of distance 2d in a straight line. It experiences air resistance\par
and rolling resistance so that the total resistance to motion when it is moving with speed v\par
\[is Av2 + R, where A and R are constants.\]
The car starts from rest and moves with constant acceleration a for a distance d. Show that\par
the work done by the engine for this half of the journey is\par
d\par
\[(ma + R + Av2) dx\]
(\par
and that it can be written in the form\par
\[w (ma + R + Av2)v\]
dv ,\par
( a\par
\[where w = \sqrt\{2ad\} .\]
For the second half of the journey, the acceleration of the car is a.\par
-\par
\textbf{(i)} In the case R > ma, show that the work done by the engine for the whole journey is\par
\[2Aad2 + 2Rd .\]
\textbf{(ii)} In the case ma 2Aad < R < ma, show that at a certain speed the driving force\par
-\par
required to maintain the constant acceleration falls to zero.\par
Thereafter, the engine does no work (and the driver applies the brakes to maintain the\par
constant acceleration). Show that the work done by the engine for the whole journey is\par
(ma R)2\par
\[2Aad2 + 2Rd + - .\]
4Aa\par
\bigskip\textbf{Question 11.} Two thin vertical parallel walls, each of height 2a, stand a distance a apart on horizontal\par
ground. The projectiles in this question move in a plane perpendicular to the walls.\par
\textbf{(i)} A particle is projected with speed \sqrt5ag towards the two walls from a point A at\par
ground level. It just clears the first wall. By considering the energy of the particle,\par
find its speed when it passes over the first wall.\par
Given that it just clears the second wall, show that the angle its trajectory makes with\par
\textasciicircum{}\circ\par
the horizontal when it passes over the first wall is 45 .\par
Find the distance of A from the foot of the first wall.\par
\textbf{(ii)} A second particle is projected with speed \sqrt5ag from a point B at ground level towards\par
\[the two walls. It passes a distance h above the first wall, where h > 0. Show that it\]
does not clear the second wall.\par
\[\geq\geq\theta\mu\nu) (\]
\end{flushleft}
\subsection*{Page 7}
\begin{flushleft}
\bigskip\textbf{Section C: Probability and Statistics}\par
12 Adam and Eve are catching fish. The number of fish, X, that Adam catches in a fixed time\par
interval T has a Poisson distribution with parameter λ. The number of fish, Y , that Eve\par
catches in the same time interval has a Poisson distribution with parameter μ. The two\par
Poisson variables are independent.\par
\textbf{(i)} By considering P(X + Y = r), show that the total number of fish caught by Adam and\par
Eve in time T also has a Poisson distribution.\par
\textbf{(ii)} Given that Adam and Eve catch a total of k fish in time T , where k is fixed, show that\par
the number caught by Adam has a binomial distribution.\par
\textbf{(iii)} Given that Adam and Eve start fishing at the same time, find the probability that the\par
first fish is caught by Adam.\par
\textbf{(iv)} You are now given that, for a Poisson distribution with parameter θ, the expected time\par
from any starting point until the next event is θ\par
-1.\par
Find the expected time from the moment Adam and Eve start fishing until they have\par
each caught at least one fish.\par
\bigskip\textbf{Question 13.} In a television game show, a contestant has to open a door using a key. The contestant is\par
\[given a bag containing n keys, where n \leq 2. Only one key in the bag will open the door.\]
There are three versions of the game. In each version, the contestant starts by choosing a\par
key at random from the bag.\par
\textbf{(i)} In version 1, after each failed attempt at opening the door the key that has been tried\par
is put back into the bag and the contestant again selects a key at random from the bag.\par
By considering the binomial expansion of (1 q)\par
-2,\par
or otherwise, find the expected\par
-\par
number of attempts required to open the door.\par
\textbf{(ii)} In version 2, after each failed attempt at opening the door the key that has been tried\par
is put aside and the contestant selects another key at random from the bag. Find the\par
expected number of attempts required to open the door.\par
\textbf{(iii)} In version 3, after each failed attempt at opening the door the key that has been tried\par
is put back into the bag and another incorrect key is added to the bag. The contes-\par
tant then selects a key at random from the bag. Show that the probability that the\par
contestant draws the correct key at the kth attempt is\par
n 1\par
- .\par
(n + k 1)(n + k 2)\par
- -\par
Show also, using partial fractions, that the expected number of attempts required to\par
open the door is infinite.\par
N\par
You may use without proof the result that as N .\par
m \to \infty \to \infty\par
\[m)=1\]
\[\geq\geq\theta\mu\nu) )\]
\end{flushleft}
\end{document}
